% appendix-b-glossary.tex — OpenBench Software Design Document, short edition (condensed from sections/appendix-b-glossary.tex on 2026-10-05).
% The domain vocabulary is generated by the Makefile (target glossary) from CONTEXT.md with the SRS generator; edit CONTEXT.md, not this file.
\section{Glossary}\label{app:glossary}

The domain vocabulary below is generated at build time from the project's shared vocabulary file and is identical to the glossary of the SRS; the design terms after it are specific to this document.

\begin{multicols}{2}
\footnotesize
\input{build/glossary-compact}

\glossgroup{Design terms}\sloppy
\glossterm{Module}{A unit of the design with one area of responsibility: one Java package under \code{edu.cwru.openbench}, or the web client. A back-end module owns its tables and exposes only its application services and the DTOs they take and return.}{}
\glossterm{Application service}{A Spring bean that is a module's public entry point, such as \code{RequestLifecycleService}. Each method runs in one transaction unless marked read-only and carries the role check for its action.}{}
\glossterm{Port}{An interface declared by a lower module and implemented by a higher one, so the lower module can ask a question without depending upward. The only port is \code{ReductionGuard}, declared by Inventory and implemented by Allocation.}{}
\glossterm{Capability interface}{An optional interface an allocator implements beside \code{Allocator} when its Allocation Class has that behavior: \code{Waitlistable}, \code{Bumpable}, or \code{ApprovalDriven}. Requests asks the \code{AllocatorRegistry} whether a class has one.}{}
\glossterm{Exclusion constraint}{A PostgreSQL constraint that rejects a row whose values conflict with an existing row under the given operators. It keeps Held or Allocated occurrences of one Resource from overlapping and Loans of one Person from overlapping.}{}
\end{multicols}
